% =====================================================================
% config/metadata.tex
% PT: Dados da dissertação. É aqui que preenche o título, o seu nome, o
%     curso, os orientadores, o júri e as datas. A capa, a folha do júri e
%     as propriedades do PDF são geradas a partir destes valores.
%     Substitua apenas o texto entre as chavetas finais { }.
% EN: Dissertation data. This is where you fill in the title, your name,
%     the programme, the supervisors, the jury and the dates. The cover,
%     the jury page and the PDF properties are generated from these values.
%     Replace only the text inside the last pair of braces { }.
% =====================================================================

% ---------------------------------------------------------------------
% Trabalho / Work
% ---------------------------------------------------------------------
% PT: Título da dissertação. Na capa é escrito em maiúsculas.
%     Ex.: Deteção de Anomalias em Redes Industriais com Aprendizagem Automática
% EN: Dissertation title. It is printed in capitals on the cover.
%     E.g.: Anomaly Detection in Industrial Networks with Machine Learning
\newcommand{\thesistitle}{Título da Dissertação}

% PT: Subtítulo opcional. Se não estiver vazio, aparece na capa e na folha
%     do júri a seguir ao título, separado por dois pontos. Deixe {} vazio
%     se não tiver subtítulo.
%     Ex.: Um Estudo de Caso no Setor Automóvel
% EN: Optional subtitle. If not empty, it is printed on the cover and the
%     jury page after the title, separated by a colon. Leave {} empty if
%     there is no subtitle.
%     E.g.: A Case Study in the Automotive Sector
\newcommand{\thesissubtitle}{}

% PT: Palavras-chave para as propriedades do PDF (separadas por vírgulas).
%     As palavras-chave visíveis estão em abstract.tex e resumo.tex.
%     Ex.: Machine Learning, Anomaly Detection, Industrial Networks
% EN: Keywords for the PDF properties (comma-separated). The visible
%     keywords are in abstract.tex and resumo.tex.
%     E.g.: Machine Learning, Anomaly Detection, Industrial Networks
\newcommand{\thesiskeywords}{Palavra-chave 1, Palavra-chave 2, Palavra-chave 3}

% ---------------------------------------------------------------------
% Aluno e curso / Student and programme
% ---------------------------------------------------------------------
% PT: Nome completo do aluno, como deve aparecer na capa (em maiúsculas).
%     Ex.: Ana Exemplo Ferreira
% EN: Student's full name, as it should appear on the cover (in capitals).
%     E.g.: Ana Exemplo Ferreira
\newcommand{\authorname}{Nome Completo do Aluno}

% PT: Nome completo do mestrado (usado nas propriedades do PDF).
%     Ex.: Mestrado em Engenharia de Software
% EN: Full name of the master's programme (used in the PDF properties).
%     E.g.: Mestrado em Engenharia de Software
\newcommand{\coursename}{Mestrado em Nome do Curso}

% PT: Primeira linha do bloco do grau na capa. Use o tipo de trabalho que
%     se aplica ao seu caso (Dissertação, Projeto ou Estágio).
%     Ex.: Relatório de Dissertação do Mestrado em
% EN: First line of the degree block on the cover. Use the type of work
%     that applies to you (Dissertação, Projeto or Estágio).
%     E.g.: Relatório de Dissertação do Mestrado em
\newcommand{\coverdegree}{Relatório de Dissertação do Mestrado em}

% PT: Segunda linha do bloco do grau na capa: a área do mestrado.
%     Ex.: Engenharia de Software
% EN: Second line of the degree block on the cover: the programme field.
%     E.g.: Engenharia de Software
\newcommand{\coursefield}{Nome do Curso}

% PT: Instituição e escola (usadas nas propriedades do PDF).
% EN: University and school (used in the PDF properties).
\newcommand{\institutionname}{Universidade Politécnica de Setúbal}
\newcommand{\schoolname}{Escola Superior de Tecnologia de Setúbal}

% PT: Cidade (disponível para uso no texto, ex.: nos agradecimentos).
% EN: City (available for use in the text, e.g. in the acknowledgements).
\newcommand{\cityname}{Setúbal}

% ---------------------------------------------------------------------
% Orientação / Supervision
% ---------------------------------------------------------------------
% PT: Orientador, com grau académico.
%     Ex.: Prof. Doutor João Exemplo Almeida
% EN: Supervisor, with academic title.
%     E.g.: Prof. Doutor João Exemplo Almeida
\newcommand{\supervisorname}{Prof. Doutor Nome do Orientador}

% PT: Coorientador (OPCIONAL). Se deixar {} vazio, a capa não mostra o
%     bloco COORIENTADOR.
%     Ex.: Prof.ª Doutora Maria Exemplo Sousa   ou   {} (sem coorientador)
% EN: Co-supervisor (OPTIONAL). If you leave {} empty, the cover does not
%     show the COORIENTADOR block.
%     E.g.: Prof.ª Doutora Maria Exemplo Sousa   or   {} (no co-supervisor)
\newcommand{\cosupervisorname}{Prof. Doutor Nome do Coorientador}

% PT: Rótulos da capa e da folha do júri. Mude para a forma feminina
%     (ORIENTADORA, COORIENTADORA, Orientadora) quando se aplicar.
% EN: Labels on the cover and the jury page. Change them to the feminine
%     form (ORIENTADORA, COORIENTADORA, Orientadora) when it applies.
\newcommand{\supervisorlabel}{ORIENTADOR}
\newcommand{\cosupervisorlabel}{COORIENTADOR}
\newcommand{\juryadvisorlabel}{Orientador}

% ---------------------------------------------------------------------
% Datas / Dates
% ---------------------------------------------------------------------
% PT: Mês (em minúsculas, em português) e ano de entrega.
%     Ex.: novembro  /  2026
% EN: Month (lowercase, in Portuguese) and year of submission.
%     E.g.: novembro  /  2026
\newcommand{\submissionmonth}{mês}
\newcommand{\submissionyear}{AAAA}

% PT: Mês e ano da prova pública (defesa). Aparecem no fundo da capa e da
%     folha do júri. Por omissão são iguais aos da entrega; depois da
%     defesa, substitua por valores concretos.
%     Ex.: \newcommand{\defencemonth}{janeiro}  /  \newcommand{\defenceyear}{2027}
% EN: Month and year of the public defence. They are printed at the bottom
%     of the cover and the jury page. By default they match the submission
%     date; after the defence, replace them with the actual values.
%     E.g.: \newcommand{\defencemonth}{janeiro}  /  \newcommand{\defenceyear}{2027}
\newcommand{\defencemonth}{\submissionmonth}
\newcommand{\defenceyear}{\submissionyear}

% ---------------------------------------------------------------------
% Júri (folha do júri) / Jury (jury page)
% ---------------------------------------------------------------------
% PT: Formato de cada membro: Grau Nome, Instituição.
% EN: Format of each member: Title Name, Institution.

% PT: Presidente do júri.
%     Ex.: Prof. Doutor Rui Exemplo Lopes, Universidade Politécnica de Setúbal
% EN: Jury president.
%     E.g.: Prof. Doutor Rui Exemplo Lopes, Universidade Politécnica de Setúbal
\newcommand{\jurypresident}{Prof. Doutor Nome do Presidente, Universidade Politécnica de Setúbal}

% PT: Orientador no júri (normalmente igual a \supervisorname + instituição).
% EN: Supervisor on the jury (usually \supervisorname + institution).
\newcommand{\juryadvisor}{Prof. Doutor Nome do Orientador, Universidade Politécnica de Setúbal}

% PT: Vogal (arguente). É obrigatório pelo menos um.
%     Ex.: Prof.ª Doutora Carla Exemplo Pinto, Universidade de Lisboa
% EN: Examiner ("vogal"). At least one is required.
%     E.g.: Prof.ª Doutora Carla Exemplo Pinto, Universidade de Lisboa
\newcommand{\juryvogal}{Prof. Doutor Nome do Vogal, Instituição}

% PT: Segundo vogal (OPCIONAL). Deixe {} vazio se o júri tiver só um
%     vogal; preencha-o se tiver dois. Cada vogal ocupa uma linha.
%     Ex.: Prof. Doutor Tiago Exemplo Nunes, Universidade de Évora
% EN: Second examiner (OPTIONAL). Leave {} empty if the jury has a single
%     examiner; fill it in if there are two. Each examiner takes one line.
%     E.g.: Prof. Doutor Tiago Exemplo Nunes, Universidade de Évora
\newcommand{\jurysecondvogal}{}
